\section{Ext groups}
\label{section-ext}

\noindent
In this section we do a tiny bit of homological algebra,
in order to establish some fundamental properties of
depth over Noetherian local rings.

\begin{lemma}
\label{lemma-resolution-by-finite-free}
Let $R$ be a ring. Let $M$ be an $R$-module.
\begin{enumerate}
\item There exists an exact complex
$$
\ldots \to F_2 \to F_1 \to F_0 \to M \to 0.
$$
with $F_i$ free $R$-modules.
\item If $R$ is Noetherian and $M$ finite over $R$, then we
can choose the complex such that $F_i$ is finite free.
In other words, we can find an exact complex
$$
\ldots \to R^{\oplus n_2} \to R^{\oplus n_1} \to R^{\oplus n_0} \to M \to 0.
$$
\end{enumerate}
\end{lemma}

\begin{proof}
Let us explain only the Noetherian case.
As a first step choose a surjection $R^{n_0} \to M$.
Then having constructed an exact complex of length
$e$ we simply choose a surjection $R^{n_{e + 1}} \to
\Ker(R^{n_e} \to R^{n_{e-1}})$ which is possible
because $R$ is Noetherian.
\end{proof}

\begin{definition}
\label{definition-finite-free-resolution}
Let $R$ be a ring. Let $M$ be an $R$-module.
\begin{enumerate}
\item A (left) {\it resolution} $F_\bullet \to M$ of $M$ is an exact complex
$$
\ldots \to F_2 \to F_1 \to F_0 \to M \to 0
$$
of $R$-modules.
\item A {\it resolution of $M$ by free $R$-modules} is a resolution
$F_\bullet \to M$ where each $F_i$ is a free $R$-module.
\item A {\it resolution of $M$ by finite free $R$-modules} is a resolution
$F_\bullet \to M$ where each $F_i$ is a finite free $R$-module.
\end{enumerate}
\end{definition}

\noindent
We often use the notation $F_{\bullet}$ to denote a complex
of $R$-modules
$$
\ldots \to F_i \to F_{i-1} \to \ldots
$$
In this case we often use $d_i$ or $d_{F, i}$ to denote the map
$F_i \to F_{i-1}$. In this section we are always going to
assume that $F_0$ is the last nonzero term in the complex.
The {\it $i$th homology group of the complex} $F_{\bullet}$
is the group $H_i = \Ker(d_{F, i})/\Im(d_{F, i + 1})$.
A {\it map of complexes $\alpha : F_{\bullet} \to G_{\bullet}$}
is given by maps $\alpha_i : F_i \to G_i$ such that
$\alpha_{i-1} \circ d_{F, i} = d_{G, i-1} \circ \alpha_i$.
Such a map induces a map on homology $H_i(\alpha) :
H_i(F_{\bullet}) \to H_i(G_{\bullet})$. If $\alpha, \beta
:  F_{\bullet} \to G_{\bullet}$ are maps of complexes, then
a {\it homotopy} between $\alpha$ and $\beta$ is given by
a collection of maps $h_i : F_i \to G_{i + 1}$ such that
$\alpha_i - \beta_i = d_{G, i + 1} \circ h_i +
h_{i-1} \circ d_{F, i}$.
Two maps $\alpha, \beta : F_{\bullet} \to G_{\bullet}$ are
said to be {\it homotopic} if a homotopy between $\alpha$
and $\beta$ exists.

\medskip\noindent
We will use a very similar notation regarding complexes
of the form $F^{\bullet}$ which look like
$$
\ldots \to F^i \xrightarrow{d^i} F^{i + 1} \to \ldots
$$
There are maps of complexes, homotopies, etc.
In this case we set $H^i(F^{\bullet}) =
\Ker(d^i)/\Im(d^{i - 1})$ and we call it
the {\it $i$th cohomology group}.

\begin{lemma}
\label{lemma-homotopic-equal-homology}
Any two homotopic maps of complexes induce the same maps on
(co)homology groups.
\end{lemma}

\begin{proof}
Omitted.
\end{proof}

\begin{lemma}
\label{lemma-compare-resolutions}
Let $R$ be a ring. Let $M \to N$ be a map of $R$-modules.
Let $N_\bullet \to N$ be an arbitrary resolution.
Let
$$
\ldots \to F_2 \to F_1 \to F_0 \to M
$$
be a complex of $R$-modules where each $F_i$ is a free $R$-module. Then
\begin{enumerate}
\item there exists a map of complexes $F_\bullet \to N_\bullet$ such that
$$
\xymatrix{
F_0 \ar[r] \ar[d] & M \ar[d] \\
N_0 \ar[r] & N
}
$$
is commutative, and
\item any two maps $\alpha, \beta : F_\bullet \to N_\bullet$ as in (1)
are homotopic.
\end{enumerate}
\end{lemma}

\begin{proof}
Proof of (1). Because $F_0$ is free we can find a map $F_0 \to N_0$
lifting the map $F_0 \to M \to N$. We obtain an induced
map $F_1 \to F_0 \to N_0$ which ends up in the image
of $N_1 \to N_0$. Since $F_1$ is free we may lift this
to a map $F_1 \to N_1$. This in turn induces a map
$F_2 \to F_1 \to N_1$ which maps to zero into
$N_0$. Since $N_\bullet$ is exact we see that
the image of this map is contained in the image
of $N_2 \to N_1$. Hence we may lift to get a map
$F_2 \to N_2$. Repeat.

\medskip\noindent
Proof of (2). To show that $\alpha, \beta$ are homotopic it suffices
to show the difference $\gamma = \alpha - \beta$ is homotopic
to zero. Note that the image of $\gamma_0 : F_0 \to N_0$
is contained in the image of $N_1 \to N_0$. Hence we may lift
$\gamma_0$ to a map $h_0 : F_0 \to N_1$. Consider the map
$\gamma_1' = \gamma_1 - h_0 \circ d_{F, 1}$. By our choice of $h_0$
we see that the image of $\gamma_1'$ is contained in
the kernel of $N_1 \to N_0$. Since $N_\bullet$ is exact
we may lift $\gamma_1'$ to a map $h_1 : F_1 \to N_2$.
At this point we have $\gamma_1 = h_0 \circ d_{F, 1}
+ d_{N, 2} \circ h_1$. Repeat.
\end{proof}

\noindent
At this point we are ready to define the groups
$\Ext^i_R(M, N)$. Namely, choose a resolution
$F_{\bullet}$ of $M$ by free $R$-modules, see Lemma
\ref{lemma-resolution-by-finite-free}. Consider
the (cohomological) complex
$$
\Hom_R(F_\bullet, N) :
\Hom_R(F_0, N) \to
\Hom_R(F_1, N) \to
\Hom_R(F_2, N) \to \ldots
$$
We define $\Ext^i_R(M, N)$ for $i \geq 0$ to be the $i$th
cohomology group of this complex\footnote{At this point
it would perhaps be more appropriate to say ``an'' in stead
of ``the'' Ext-group.}. For $i < 0$ we set $\Ext^i_R(M, N) = 0$.
Before we continue we point out that
$$
\Ext^0_R(M, N) = \Ker(\Hom_R(F_0, N) \to \Hom_R(F_1, N)) =
\Hom_R(M, N)
$$
because we can apply part (1) of Lemma \ref{lemma-hom-exact} to
the exact sequence $F_1 \to F_0 \to M \to 0$.
The following lemma explains
in what sense this is well defined.

\begin{lemma}
\label{lemma-ext-welldefined}
Let $R$ be a ring. Let $M_1, M_2, N$ be $R$-modules.
Suppose that $F_{\bullet}$ is a free resolution of the module $M_1$,
and $G_{\bullet}$ is a free resolution of the module $M_2$.
Let $\varphi : M_1 \to M_2$ be a module map.
Let $\alpha : F_{\bullet} \to G_{\bullet}$ be
a map of complexes inducing $\varphi$ on
$M_1 = \Coker(d_{F, 1}) \to M_2 = \Coker(d_{G, 1})$,
see Lemma \ref{lemma-compare-resolutions}.
Then the induced maps
$$
H^i(\alpha) :
H^i(\Hom_R(F_{\bullet}, N))
\longrightarrow
H^i(\Hom_R(G_{\bullet}, N))
$$
are independent of the choice of $\alpha$.
If $\varphi$ is an isomorphism, so are all the maps
$H^i(\alpha)$. If $M_1 = M_2$, $F_\bullet = G_\bullet$, and
$\varphi$ is the identity, so are all the maps $H_i(\alpha)$.
\end{lemma}

\begin{proof}
Another map $\beta : F_{\bullet} \to G_{\bullet}$
inducing $\varphi$ is homotopic to $\alpha$ by
Lemma \ref{lemma-compare-resolutions}. Hence the
maps $\Hom_R(F_\bullet, N) \to
\Hom_R(G_\bullet, N)$ are homotopic.
Hence the independence result follows from
Lemma \ref{lemma-homotopic-equal-homology}.

\medskip\noindent
Suppose that $\varphi$ is an isomorphism.
Let $\psi : M_2 \to M_1$ be an inverse.
Choose $\beta : G_{\bullet} \to F_{\bullet}$
be a map inducing $\psi :
M_2 = \Coker(d_{G, 1}) \to M_1 = \Coker(d_{F, 1})$,
see Lemma \ref{lemma-compare-resolutions}.
OK, and now consider the map
$H^i(\alpha) \circ H^i(\beta) =
H^i(\alpha \circ \beta)$. By the above the
map $H^i(\alpha \circ \beta)$ is the {\it same}
as the map $H^i(\text{id}_{G_{\bullet}}) = \text{id}$.
Similarly for the composition $H^i(\beta) \circ H^i(\alpha)$.
Hence $H^i(\alpha)$ and $H^i(\beta)$ are inverses of each other.
\end{proof}

\begin{lemma}
\label{lemma-long-exact-seq-ext}
Let $R$ be a ring. Let $M$ be an $R$-module.
Let $0 \to N' \to N \to N'' \to 0$ be a
short exact sequence. Then we get a long exact
sequence
$$
\begin{matrix}
0
\to \Hom_R(M, N')
\to \Hom_R(M, N)
\to \Hom_R(M, N'')
\\
\phantom{0\ }
\to \Ext^1_R(M, N')
\to \Ext^1_R(M, N)
\to \Ext^1_R(M, N'')
\to \ldots
\end{matrix}
$$
\end{lemma}

\begin{proof}
Pick a free resolution $F_{\bullet} \to M$.
Since each of the $F_i$ are free we see that
we get a short exact sequence of complexes
$$
0 \to
\Hom_R(F_{\bullet}, N') \to
\Hom_R(F_{\bullet}, N) \to
\Hom_R(F_{\bullet}, N'') \to
0
$$
Thus we get the long exact sequence from
the snake lemma applied to this.
\end{proof}

\begin{lemma}
\label{lemma-reverse-long-exact-seq-ext}
Let $R$ be a ring. Let $N$ be an $R$-module.
Let $0 \to M' \to M \to M'' \to 0$ be a
short exact sequence. Then we get a long exact
sequence
$$
\begin{matrix}
0
\to \Hom_R(M'', N)
\to \Hom_R(M, N)
\to \Hom_R(M', N)
\\
\phantom{0\ }
\to \Ext^1_R(M'', N)
\to \Ext^1_R(M, N)
\to \Ext^1_R(M', N)
\to \ldots
\end{matrix}
$$
\end{lemma}

\begin{proof}
Pick sets of generators $\{m'_{i'}\}_{i' \in I'}$ and
$\{m''_{i''}\}_{i'' \in I''}$ of $M'$ and $M''$.
For each $i'' \in I''$ choose a lift $\tilde m''_{i''} \in M$
of the element $m''_{i''} \in M''$. Set $F' = \bigoplus_{i' \in I'} R$,
$F'' = \bigoplus_{i'' \in I''} R$ and $F = F' \oplus F''$.
Mapping the generators of these free modules to the corresponding
chosen generators gives surjective $R$-module maps $F' \to M'$,
$F'' \to M''$, and $F \to M$. We obtain a map of short exact sequences
$$
\begin{matrix}
0 & \to & M' & \to & M & \to & M'' & \to & 0 \\
& & \uparrow & & \uparrow & & \uparrow \\
0 & \to & F' & \to & F & \to & F'' & \to & 0 \\
\end{matrix}
$$
By the snake lemma we see that the sequence of kernels
$0 \to K' \to K \to K'' \to 0$ is short exact sequence of $R$-modules.
Hence we can continue this process indefinitely. In other words
we obtain a short exact sequence of resolutions fitting into the diagram
$$
\begin{matrix}
0 & \to & M' & \to & M & \to & M'' & \to & 0 \\
& & \uparrow & & \uparrow & & \uparrow \\
0 & \to & F_\bullet' & \to & F_\bullet & \to & F_\bullet'' & \to & 0 \\
\end{matrix}
$$
Because each of the sequences $0 \to F'_n \to F_n \to F''_n \to 0$
is split exact (by construction) we obtain a short exact sequence of
complexes
$$
0 \to
\Hom_R(F''_{\bullet}, N) \to
\Hom_R(F_{\bullet}, N) \to
\Hom_R(F'_{\bullet}, N) \to
0
$$
by applying the $\Hom_R(-, N)$ functor.
Thus we get the long exact sequence from
the snake lemma applied to this.
\end{proof}

\begin{lemma}
\label{lemma-annihilate-ext}
Let $R$ be a ring. Let $M$, $N$ be $R$-modules.
Any $x\in R$ such that either $xN = 0$, or $xM = 0$
annihilates each of the modules $\Ext^i_R(M, N)$.
\end{lemma}

\begin{proof}
Pick a free resolution $F_{\bullet}$ of $M$.
Since $\Ext^i_R(M, N)$
is defined as the cohomology of the complex
$\Hom_R(F_{\bullet}, N)$ the lemma is
clear when $xN = 0$. If $xM = 0$, then
we see that multiplication by $x$ on $F_{\bullet}$
lifts the zero map on $M$. Hence by Lemma
\ref{lemma-ext-welldefined} we see that it
induces the same map on Ext groups as the
zero map.
\end{proof}

\begin{lemma}
\label{lemma-ext-noetherian}
Let $R$ be a Noetherian ring. Let $M$, $N$ be finite $R$-modules.
Then $\Ext^i_R(M, N)$ is a finite $R$-module for all $i$.
\end{lemma}

\begin{proof}
This holds because $\Ext^i_R(M, N)$ is computed as the
cohomology groups of a complex $\Hom_R(F_\bullet, N)$
with each $F_n$ a finite free $R$-module, see
Lemma \ref{lemma-resolution-by-finite-free}.
\end{proof}




